\documentclass[11pt]{article}
\usepackage{fontspec}
\usepackage[french,english]{babel}
\usepackage{amsmath}
\usepackage{amssymb}
\usepackage{amsthm}
\usepackage{mathrsfs}
\usepackage{tikz-cd}
\usepackage{enumitem}

\newcommand{\printedpage}[1]{\marginpar{\footnotesize p.~#1}}
\newcommand{\msp}[1]{\marginpar{\footnotesize [#1]}}
\newcommand{\Fib}{\underline{\mathrm{Fib}}}
\newcommand{\Ch}{\mathrm{Ch}}
\newcommand{\pr}{\mathrm{pr}}
\newcommand{\circled}[1]{\tikz[baseline=(c.base)]\node[draw,circle,inner sep=1pt](c){#1};}

\begin{document}

\printedpage{64}
% head: A. Grothendieck : Dérivateurs, Ch. VI, Édition des Universités de Montpellier II et Paris VII

\[
\begin{tikzcd}[column sep=large]
 & \Ch(S) \arrow[d, "p_S"] \arrow[dl, dotted, "\beta_S"'] \arrow[dd, dotted, bend right=40, "\sigma_S"'] \arrow[dr, phantom, "\text{cart.}"]
 & C_S(X)=C(X) \arrow[l, "\pi"'] \arrow[d, "p'"] \arrow[dll, bend right=40, "\bar f_X=f_{C(X)}"'] \\
S & S\times S \arrow[l, "\pr_2^S"'] \arrow[d, "\pr_1^S"] \arrow[dr, phantom, "\text{cart.}"]
 & X\times S \arrow[l, "\pi_1=f_X\times\mathrm{id}_S"] \arrow[d, "\pr_1"] \\
 & S & X(=S'). \arrow[l, "f_X"'] \arrow[uu, dotted, bend right=40, "i_X"']
\end{tikzcd}
\]

It is the inverse image of $(\Ch(S),\pr_1^Sp_S=\sigma_S)$ under $X\xrightarrow{f_X}S$, but we have made the factorization of $\sigma_S$ explicit in the diagram to reinterpret the structural arrow $\bar f_X=f_{C(X)}$, defined as
\[
f_{C(X)}=\beta_S\pi(=\pr_2^S\underbrace{p_S\pi}_{=\pi_1p'})
\]
by the formula
\[
f_{C(X)}=\underbrace{(\pr_2^S\pi_1)}_{\text{projection } \pr_2 \text{ from } X\times S \text{ to } S}p'.
\]

\msp{page 102}
Now
\[
\begin{aligned}
&p'\in\mathrm{Fib}, && \text{since obtained from } p_S\in\mathrm{Fib} \text{ by base change } \pi_1,\\
&(\pr_2 : X\times S\longrightarrow S)\in\mathrm{Fib} && \text{provided that } X\longrightarrow\underbrace{e}_{\text{final object of } \mathcal{M}}\quad \text{is in Fib (same reason),}
\end{aligned}
\]
so if $X\in\Fib\,e$, we have $f_{C(X)}\in\mathrm{Fib}$. Since this is an essential property that we want to have for \textit{every} $X$ in $\mathcal{M}/S$, we are going to make the additional hypothesis

\textbf{Hypothesis 3.} \textit{$\forall\,X\in\mathit{Ob}\,\mathcal{M}$, $X\longrightarrow e$ is in Fib.}

As for
\[
i_X : X\longrightarrow C(X),
\]
it is obtained from the section $\alpha_S$ of $(\Ch(S),\sigma_S)$ over $S$ by base change $X\longrightarrow S$. Let us now note that
\[
\begin{tikzcd}
\Ch(S) \arrow[r, shift left, "\sigma_S"] \arrow[r, shift right, "\beta_S"'] & S
\end{tikzcd}
\qquad\qquad\qquad \text{are in } W,
\]
since their composites with $S\xrightarrow{\alpha_S}\Ch(S)$ are the identity, and $\alpha_S\in W$. On the other hand, $\sigma_S=\pr_1^S\circ p$ is in Fib, as a composite of two elements in Fib (NB $S\longrightarrow e$ in Fib by Hypothesis 3). Thus $\sigma_S$, $\beta_S$ are in $\mathrm{Fib}\cap W$. We shall need % [sic] \pr_1^S\circ p -> \pr_1^S\circ p_S, la flèche est notée p_S dans le diagramme et plus haut (\pr_1^Sp_S=\sigma_S)

\textbf{Hypothesis 4.} \textit{The set of arrows $\mathit{Fib}\cap W$ is stable under every base change.}

% foot: Version 9 novembre 2001

\printedpage{65}
% head: A. Grothendieck : Dérivateurs, Ch. VI, Édition des Universités de Montpellier II et Paris VII

\msp{page 103}
This therefore implies that
\[
C(X)\xrightarrow{\pr_1\pi'}X
\]
\texttt{[rather} $\pr_1p'$\texttt{]} is $\in W$ (being obtained from $\sigma_S : \Ch(S)\longrightarrow S$ by base change), so every section $i_X : X\longrightarrow C(X)$ of $C(X)$ over $X$ is in $i_X$ \texttt{[rather in} $W$\texttt{]}. We thus do indeed have a factorization, functorial in $X$
\[
\begin{tikzcd}
 & C(X) \arrow[dr, "\bar f_X=f_{C(X)}"] & \\
X \arrow[ur, "i_X"] \arrow[rr, "f_X"] & & S
\end{tikzcd}
\]
with
\[
i_X\in W,\qquad\qquad f_{C(X)}\in\mathrm{Fib}.
\]
This is all we had used to establish
\[
(\Fib\,S)W_S^{-1}\xrightarrow{\sim}(\mathcal{M}/S)W_S^{-1}.
\]
But it is now a matter, under the hypothesis where we introduce an $\mathcal{L}$, of proving that for every $X$ in $\mathcal{M}/S$, the two arrows
\[
\begin{tikzcd}
C(X) \arrow[r, shift left, "C(i_X)"] \arrow[r, shift right, "i_{C(X)}"'] & C(C(X))
\end{tikzcd}
\]
are equal in $(\mathcal{M}/S)(W_S^{\mathcal{L}})^{-1}$. Since by hypothesis we have $W_S^{\mathcal{L}}=W_S$ in $\Fib\,S$ (Hypothesis 1), and $C(X),\ C(C(X))\in\mathrm{Ob}\,\Fib\,S$,
\msp{page 104}
it suffices to prove that the two arrows are equal \textit{in} $(\Fib\,S)W_S^{-1}$. This statement is independent of the datum of $\mathcal{L}$, and we feel like proving it.

Perhaps there is a way to make a geometric language explicit, to find our way around. Let us identify the objects $X$ of $\mathcal{M}$ with the functors $\mathcal{M}^o\longrightarrow(\mathrm{Ens})$ that they define,
\[
X\longmapsto(T\longmapsto X(T)=\mathrm{Hom}_{\mathcal{M}}(T,X))\,.
\]
Thus, for every $T$ % [sic] pout -> pour, faute de frappe
\[
\Ch(X)(T)
\]
is a set, called \textit{the set of paths in $X(T)$.} It maps (by $p_X$) functorially (in $T$) to $(X\times X)(T)=X(T)\times X(T)$, i.e. every path $c$ in $X(T)$ has an origin $\sigma(c)$, and a target $\beta(c)$. Moreover, by $\alpha_X$, $X(T)$ maps (functorially in $T$) to the set of paths of $X(T)$: for every element $x$ in $X(T)$, there is therefore a well-defined path (denoted $1_x$ or $\alpha_x$) with source and target $x$, which we call the trivial, or empty, path. This said,

% foot: Version 9 novembre 2001

\printedpage{66}
% head: A. Grothendieck : Dérivateurs, Ch. VI, Édition des Universités de Montpellier II et Paris VII

\msp{page 105}
for $X$ in $\mathcal{M}/S$, $C(X)$ represents the following contrafunctor in $\mathcal{M}$: for every $T$ in $\mathcal{M}$,
\[
C(X)(T)=\Bigl\{(x,c)\Bigm| x\in X(T),\ c : f_X(x)\xrightarrow{c}s\ \text{ a path in } S(T), \text{ with source } f_X(x)\Bigr\}\,,
\]
where $f_X : X\longrightarrow S$ is the structural morphism.

The ``calculations'' done, two (?) days ago show that we likewise have an interpretation
\[
\begin{aligned}
C(C(X))(T)=\Bigl\{(x,c,c')\Bigm| x\in X(T),\ f_X(x)\xrightarrow{c}s\xrightarrow{c'}t,\\
c\text{ a path with source } f_X(x),\\
c'\text{ a path with source equal to } \beta(c)\Bigr\}
\end{aligned}
\]
and the maps $i_{C(X)}$ and $C(i_X)$ are identified respectively with
\[
\begin{aligned}
(x,c)\quad&\longmapsto\quad(x,1_{f(x)},c)\quad f(x)\xrightarrow{1_x}f(x)\xrightarrow{c}s\\
(x,c)\quad&\longmapsto\quad(x,c,1_{\beta(c)})\quad f(x)\longrightarrow s\xrightarrow{1_s}s
\end{aligned}
\]
\texttt{[rather} $f(x)\xrightarrow{1_{f(x)}}f(x)\xrightarrow{c}s$\texttt{]}. To compare them, we have an irresistible urge to write
\[
c\circ 1_x=1_s\circ c=c,
\]
but do we have the right? We would have to be able to compose the paths, and that is very, very difficult!

But we can always axiomatize. We therefore suppose that for $X$ in $\mathcal{M}$, $\Ch(X)$
\msp{page 106}
is not just an amorphous structure of (so-called) ``paths'' in $X$, just the datum of the source and target of the paths and of the trivial paths (these structures on $\Ch(X)$ depending functorially on $X$), but that we have a composition of paths when they are ``composable'', i.e. an arrow from a fiber product
\[
(c,c')\longmapsto c'\circ c : (\Ch(X),\beta_X)\times_X(\Ch(X),\sigma_X)\xrightarrow{\pi_X}\Ch(X),
\]
i.e. we can compose paths in $X(T)$
\[
x\xrightarrow{c}y\xrightarrow{c'}z
\]
if and only if $\text{target}\,c=\text{source}\,c'$, and moreover we of course require that $\text{source}(c'\circ c)=\text{source}\,c$, $\text{target}(c'\circ c)=\text{target}\,c'$, which is expressed by a commutative diagram
\[
\begin{tikzcd}
(\Ch(X),\beta_X)\times_X(\Ch(X),\sigma_X) \arrow[rr] \arrow[dr] & & \Ch(X) \arrow[dl, "p_X"] \\
 & X\times X. &
\end{tikzcd}
\]

% foot: Version 9 novembre 2001

\printedpage{67}
% head: A. Grothendieck : Dérivateurs, Ch. VI, Édition des Universités de Montpellier II et Paris VII

The natural thought is that with this composition and these identities, the $X(T)$, $\Ch(T)$ become the objects and the arrows of a \textit{category} - i.e. that $(X,\Ch(X),\alpha_X,p_X,\pi_X)$ define a structure of ``category in $\mathcal{M}$''. But for what we want to do now, we probably do not need everything, i.e.
\msp{page 107}
the relevant axiom of associativity of composition and the one on units, but only that the $1_x$ be two-sided units for composition.

Supposing this additional structure on the functor Ch given, and returning to the consideration of the associated functors % [sic] supossant -> supposant, faute de frappe ; foncteurs -> foncteur, accord avec « du » et « associé »
\[
C : \mathcal{M}/S\longrightarrow\Fib\,S,
\]
we find a canonical morphism
\[
C(C(X))\xrightarrow{\gamma_X}C(X)
\]
by
\[
(x,c,c')\longmapsto(x,c'\circ c),
\]
and the axiom of two-sided units assures us at once that
\[
(\ast)\qquad\qquad\gamma_X\circ i_{C(X)}=\gamma_X\circ C(i_X)=\mathrm{id}_{C(X)}.
\]
Since $i_{C(X)}\in W$, this already proves that $\gamma_X\in W$, so $\gamma_X\in W_S$. So $\gamma_X$ becomes an isomorphism in $(\Fib\,S)W_S^{-1}$, and the relation $(\ast)$ does indeed imply $i_{C(X)}=C(i_X)$ in the localized category, q.e.d.

I am going to summarize again what I have worked out.

\msp{page 108}
\textbf{Theorem-Scholie 2.} Let $\mathcal{M}$ be a category, equipped with two sets of arrows
\[
W,\ \mathrm{Fib}\subseteq\mathrm{Fl}(\mathcal{M}).
\]
We make the following hypotheses.

\begin{enumerate}[label={}]
\item[$1^o$] $W$ contains the isomorphisms, if $f$, $g$ are composable, if two of $f$, $g$, $gf$ are in $W$, the third also.
\item[$2^o$] $\mathcal{M}$ is stable under finite $\varprojlim$, Fib contains the isomorphisms, is stable under composition and under base change (\footnote[59]{These stability hypotheses on Fib will be used only in II, and even then only in case A), taken up again in toto in the form A$'$ below (p.~112).}).
\end{enumerate}

For $S$ in $\mathcal{M}$, \fbox{$\Fib\,S$} denotes the full subcategory of $\mathcal{M}/S$ formed by the $X\longrightarrow S$ which are in Fib. We distinguish two different versions of an axiom of factorization of an arbitrary $f : X\longrightarrow S$ in $\mathcal{M}$ as $X\xrightarrow{i}\bar X\xrightarrow{\bar f}S$, with $i\in W$, $\bar f\in\mathrm{Fib}$.

% foot: Version 9 novembre 2001

\printedpage{68}
% head: A. Grothendieck : Dérivateurs, Ch. VI, Édition des Universités de Montpellier II et Paris VII

\begin{enumerate}[label={}]
\item[A)] We are given, for every $S$ in $\mathcal{M}$, a functor
\[
\mathcal{M}/S\longrightarrow\mathcal{M}/S\qquad\qquad X\longmapsto C(X)\ \ (=\bar X)
\]
and a morphism of functors in $\mathrm{X}$
\[
X\xrightarrow{i_X}C(X)
\]
with
\[
i_X\in W,\qquad\qquad C(X)\in\mathrm{Ob}\,\Fib\,S.
\]
\item[B)] Denoting by
\[
\mathrm{Cof}_0\subseteq\mathrm{Fl}\,\mathcal{M}
\]
the set of arrows of $\mathcal{M}$ which have the right lifting property \texttt{[rather left]} with respect to Fib, we suppose ``in the style of \textsc{Quillen}'', that every arrow $f : X\longrightarrow S$ factors as
\[
X\xrightarrow{i_X}C(X)\xrightarrow{\bar f}S,
\]
with
\[
i\in\mathrm{Cofib}_0\cap W,\qquad\qquad\bar f\in\mathrm{Fib}.
\]
\texttt{[where} $\mathrm{Cof}_0=\mathrm{Cofib}_0$\texttt{]}.
\end{enumerate}

\msp{page 109}
Let us take $S$ in $\mathcal{M}$, consider
\[
W_S\subseteq\mathrm{Fl}(\mathcal{M}/S)
\]
the set of arrows induced by $W$ (
\begin{tikzcd}[column sep=small, row sep=small, baseline={([yshift=.5ex]current bounding box.south)}]
f : X \arrow[rr] \arrow[dr] & & Y \arrow[dl] \\
 & S &
\end{tikzcd}
is in $W_S$ if and only if it is on forgetting $S$), and let us denote by $W_S$ (by abuse of notation) the set of arrows induced by $W_S$ on $\Fib\,S$, or more generally on any full subcategory $\mathcal{F}$ of $\mathcal{M}/S$, with
\[
\Fib\,S\subseteq\mathcal{F}\subseteq\mathcal{M}/S.
\]

\circled{I} The functor
\[
(\ast)\qquad\qquad(\Fib\,S)W_S^{-1}\xrightarrow{\sim}(\mathcal{M}/S)W_S^{-1}
\]
induced by the inclusion is an equivalence of categories, more generally for every $\mathcal{F}$ as above
\[
(\Fib\,S)W_S^{-1}\longrightarrow\mathcal{F}W_S^{-1}\longrightarrow(\mathcal{M}/S)W_S^{-1}
\]
are equivalences of categories. A quasi-inverse functor is obtained (in both cases A) and B)) with the help of the functor
\[
\gamma : \mathcal{M}/S\longrightarrow(\Fib\,S)W_S^{-1}
\]

% foot: Version 9 novembre 2001

\end{document}
